\subsection{Regular descent of the oriented ruling}

\begin{lemma}[Oriented relative Segre descent]
\label{lem:oriented-segre-descent}
Let \(B\) be a connected projective complex variety with
\(H^0(B,\OO_B)=\C\), and let \(V,V'\) have the same dimension
\(n\ge2\).  Let \(\mathcal U,\mathcal U'\) be vector bundles
of rank \(n\) on \(B\).  Suppose a linear bundle isomorphism
\[
 \alpha:V\otimes\mathcal U^*\xrightarrow{\sim}
                         V'\otimes\mathcal U'^*
\]
preserves the rank-one cones and the ruling whose parameter bundle
is \(\Pj(V)\times B\).  The induced isomorphism
\(\Pj(V)\times B\simeq\Pj(V')\times B\) is regular and is
induced by one constant projective isomorphism \([g]:\Pj(V)\to\Pj(V')\).
Locally \(\alpha(T)=gTh^{-1}\), with the same \([g]\) on every
chart.  The possible line-bundle twists of vector-bundle lifts do
not change the projective left coefficient lines of any homogeneous
functorial coefficient construction.
\end{lemma}
\begin{proof}
Consider the relative Grassmannian of rank-\(n\) linear subspaces
of \(V\otimes\mathcal U^*\).  Requiring the universal subspace
to lie in the rank-one cone is a closed condition, given by the
restrictions of its quadratic minors.  Its relative parameter
scheme is the disjoint union
\[
 (\Pj(V)\times B)\ \sqcup\ \Pj(\mathcal U^*).
\]
This is a statement about schemes, not just their points.  It can
be checked on a bundle frame.  Every maximal rank-one subspace
fixes one factor.  At \(v\otimes U^*\), the linearized minor
equations force a deformation to be
\(u\mapsto \dot v\otimes u\) modulo \(v\otimes U^*\);
its tangent space is \(\Hom(\C v,V/\C v)\).
The other family has the analogous tangent space.  Each component
therefore has precisely the smooth projective-space scheme
structure furnished by its displayed parameterization.  These
identifications are natural on overlapping bundle frames.
The isomorphism \(\alpha\) induces a regular map of these relative
parameter schemes, and the orientation hypothesis preserves the
first component.

For completeness, an isomorphism of two trivial \(\Pj^{n-1}\)-bundles
over any base gives a regular projective-matrix-valued map, not
merely a fibrewise collection.  Its pullback of the relative
\(\OO(1)\) has the form \(\OO(1)\otimes p^*L\) for a line
bundle \(L\) on the base.  To see this, tensor by \(\OO(-1)\);
the resulting line bundle is trivial on every projective-space
fibre, and its pushforward and evaluation identify it with a
pullback from the base.  Pushing forward \(\OO(1)\) now gives
an isomorphism of the corresponding trivial rank-\(n\) bundles,
with that line-bundle twist.  On charts trivializing \(L\), it
is an invertible matrix \(g_i\) of regular functions.  On overlaps,
\(g_j=u_{ij}g_i\) for an invertible regular function \(u_{ij}\).
The classes \([g_i]\) thus glue to a regular morphism to the
projective-isomorphism scheme, identified after choosing fixed
bases with \(\PGL_n\).  This accounts explicitly for the possible
absence of a global linear lift.

The group \(\PGL_n\) is affine: it is the determinant-nonzero
principal open in the projective space of matrices.  A morphism
from \(B\) to an affine scheme factors through
\(\Spec H^0(B,\OO_B)=\Spec\C\), so the projective map is constant.
Choose a constant linear lift \(g\).  Each local \(g_i\) then
differs from \(g\) by a scalar unit.  After removing its two
projective-factor actions, a linear map preserving every
decomposable line is scalar: apply it to sums of two tensors with
one common factor.  Consequently the remaining action is a right
bundle map, and the scalar units may be absorbed in its local
lifts \(h\).  This proves the local tensor formula.

Finally a scalar multiplying a lift acts on a homogeneous functor
of degree \(a\) by its scalar to the power \(a\) (or \(-a\)
on its dual).  It cannot alter a projective coefficient line.
For the degree-twelve construction below, both left Cauchy factors
have the same degree twelve; the exterior-duality determinant
multiplier is common as well.  Thus neither line-bundle transition
functions nor right changes of frame can introduce two different
projective transformations on the left.
\end{proof}

The lemma applies after pulling the target bundles back along the
intrinsic base isomorphism \(\sigma:\Sigma_0(R)\to\Sigma_0(R')\).
Here the orientation is supplied by
Lemma~\ref{lem:deepest-normal-cone}; it is not an extra marking of
the failure scheme.  Nor is a canonical basis of the trivial ruling
required: different choices alter its representative only by constant
projective coordinate changes.

\subsection{The intrinsic residual ideal}

\begin{lemma}[Residual graded data without a determinant equation]
\label{lem:intrinsic-residual-ideal}
On \(B=\Sigma_0(R)\), put
\(\mathcal A=\Sym(V^*\otimes\mathcal K_0)\), and let
\(\mathcal I\subset\mathcal A\) be the homogeneous ideal of
\(\mathscr C_R\) in its intrinsic degree-one ambient bundle.
Let \(\mathcal D\subset\mathcal A\) be the ideal of the reduced
cone.  Its degree-four part is the line bundle
\begin{equation}\label{eq:determinant-line-bundle}
 L=\mathcal D_4=(\det V)^*\otimes\det\mathcal K_0
                  \ \subset\ \mathcal A_4,
 \qquad \mathcal D=L\otimes\mathcal A(-4).
\end{equation}
The graded ideal
\begin{equation}\label{eq:intrinsic-colon-readout}
 \mathcal J=(\mathcal I:\mathcal D)
   =\{a\in\mathcal A:a\mathcal D\subset\mathcal I\}
\end{equation}
is intrinsic, and its degree-twelve part is the maximal-minor
coefficient space used in the Schur readout.  Multiplication gives
canonical isomorphisms
\begin{equation}\label{eq:twisted-residual-piece}
 L\otimes\mathcal J_{12}\xrightarrow{\sim}\mathcal I_{16},
 \qquad \mathcal J_{12}\simeq L^*\otimes\mathcal I_{16}.
\end{equation}
The construction is preserved by graded cone isomorphisms.  Its
local residual-ideal calculation also commutes with arbitrary base
change when \(\mathcal D\) is transported as the determinant ideal.
\end{lemma}
\begin{proof}
The graded cone algebra determines \(\mathcal A\) and its
homogeneous kernel \(\mathcal I\).  Its reduction determines
\(\mathcal D\).  In frames for \(V\) and \(\mathcal K_0\),
Lemma~\ref{lem:deepest-normal-cone} says
\(\mathcal D=(d)\) and \(\mathcal I=dJ\), where
\(d=\det T\) and \(J=I_6(\gamma_R\Sym^2T)\).
The determinant coefficient spans
\(\bigwedge^4V^*\otimes\bigwedge^4\mathcal K_0\) inside
\(\Sym^4(V^*\otimes\mathcal K_0)\), proving
\eqref{eq:determinant-line-bundle} with its precise twist.

Multiplication by the generic determinant is injective over any
coefficient ring.  For example, choose a diagonal term order; its
leading coefficient is a unit, so the leading term of \(df\)
cannot vanish for any nonzero polynomial \(f\).  Therefore
\((dJ:d)=J\).  The ideal \(J\) is generated in degree twelve,
which gives \eqref{eq:twisted-residual-piece}.  If a frame changes
\(d\) to \(ud\), with \(u\) a unit, then
\((dJ:d)=(dJ:ud)\).  The local residual ideals hence glue in
\(\mathcal A\); the inverse line-bundle twist in
\eqref{eq:twisted-residual-piece} is exactly what cancels the
transition functions of the determinant generators.

Every graded isomorphism transports \(\mathcal I\),
\(\mathcal D\), and their colon, so it preserves this operation
before any scalar equation has been chosen.  After arbitrary base
change the same identity \(I_T=d_TJ_T\) holds for the image ideals
in the new polynomial algebra.  The determinant remains a
nonzerodivisor there, so \((I_T:d_T)=J_T\) again.  No general
claim that colons commute with base change is being used: it is the
factorization and this universal nonzerodivisor that prove it here.
On a nonreduced new base one transports the determinant ideal;
one does not recompute the absolute reduction of the total cone,
which need not commute with base change.
\end{proof}
